% Section 6: Integration with the BX3 Framework
\section{Integration with the BX3 Framework}
\label{sec:rs-integration}

The Recursive Spawning Protocol is not an isolated mechanism — it is a structural instantiation of the BX3 three-layer accountability architecture. Every RSP mechanism maps onto exactly one BX3 layer; each such mapping is minimal and structurally necessary. This section specifies the three layer mappings and demonstrates why each is exclusive to its layer.

% -------------------------------------------------------
\subsection{Purpose Layer: Accountability Anchor}
\label{sec:rs-purpose-anchor}

The Purpose Layer carries two roles in RSP that no other layer can perform.

\medskip
\noindent\textbf{Exclusive origin of spawn authorization.} At root provisioning, the Purpose Layer defines and records the spawn authorization policy $P_{\mathrm{spawn}}$ in the Fact Layer authorization ledger. $P_{\mathrm{spawn}}$ requires two conditions for any child spawn: (i) the child operating scope must be a strict subset of the parent scope ($R(n_{i+1}) \subset R(n_i)$), and (ii) the spawn must be justified by an operational necessity that cannot be met within the parent's scope. The Purpose Layer originates this policy; the Bounds Engine evaluates against it; the Fact Layer enforces it. No machine actor may issue a spawn warrant in the absence of a matching Purpose Layer authorization record.

\medskip
\noindent\textbf{Exclusive terminus of escalation.} When the chain reaches $D_{\mathrm{ESCALATE}}$, the protocol compulsory-escalates to $h(n)$ — the human anchor established at root provisioning. The human anchor is non-collapsing: for all nodes $n_i$ in chain $C$, $h(n_i) = h(n_0)$. The anchor does not migrate as the chain deepens. The human issues one of three directives — \texttt{ACK}, \texttt{RESOLVE}, or \texttt{REJECT} — and no machine actor may issue any of these. The chain terminates in human judgment, never in autonomous resolution.

The structural necessity of the exclusive terminus is immediate: if any machine actor could absorb an exception or issue a terminal directive, the accountability chain would terminate at that actor rather than at a human. The Purpose Layer's exclusivity here is not a design preference — it is the load-bearing constraint that makes the upstream BX3 guarantee structurally operative.

% -------------------------------------------------------
\subsection{Bounds Engine: Depth Enforcement}
\label{sec:rs-bounds-depth}

The Bounds Engine provides constrained execution evaluation: it evaluates proposed actions against the operating range defined by the Purpose Layer and enforces the depth constraints that prevent unbounded chain growth.

\medskip
\noindent\textbf{Depth evaluation at spawn.} At each spawn event, the Bounds Engine evaluates current chain depth $|C|$ against the Purpose-Layer-defined maximum $D_{\mathrm{max}}$, both recorded in the Fact Layer authorization ledger. If $|C| \geq D_{\mathrm{max}}$, the Bounds Engine does not execute the spawn. It transitions to \texttt{ESCALATING} state and initiates the Bailout Protocol, propagating the exception upward to $h(n)$. This enforcement is deterministic: the Bounds Engine has no discretion to exceed $D_{\mathrm{max}}$ regardless of urgency or operational exigency.

\medskip
\noindent\textbf{Pre-commit enforcement of depth and scope.} The Bounds Engine's depth evaluation is reinforced by the Fact Layer pre-commit hook. Before any spawn operation is recorded, the Fact Layer verifies: (1) the scope subset relationship $R(n_{i+1}) \subseteq R(n_i)$ holds, and (2) the resulting depth $\mathrm{depth}(n_{i+1}) \leq D_{\mathrm{max}}$ holds. The Fact Layer rejects any spawn operation violating either condition. The Bounds Engine cannot execute a rejected spawn; it may only transition to \texttt{ESCALATING}. The depth bound and scope refinement constraint are structural invariants — not policy conventions overridable in exigent circumstances.

The structural necessity of the Bounds Engine follows from the depth-limit insufficiency theorem: depth limits alone cannot prevent scope creep drift. The Bounds Engine enforces the depth constraint; the Fact Layer enforces the scope refinement constraint; together they enforce both conjuncts of the drift prevention condition.

% -------------------------------------------------------
\subsection{Fact Layer: Forensic Ledger}
\label{sec:rs-fact-ledger}

The Fact Layer provides deterministic enforcement and forensic auditability.

\medskip
\noindent\textbf{Append-only forensic ledger.} The Fact Layer records every spawn event, every state transition, every Purpose Layer directive received, and every unresolved exception. Each entry is timestamped, attributed to the node that generated it, and hash-chained to the preceding entry. Hash-chaining ensures no entry can be inserted, deleted, reordered, or altered after recording: any tampering breaks hash continuity and is detectable by any observer with ledger read access.

\medskip
\noindent\textbf{Immutable parent pointer.} Once $\alpha(n)$ is established at node provisioning, no actor — including the Purpose Layer after provisioning — may modify it. The Fact Layer write protocol rejects any attempt to overwrite $\alpha(n)$. Chain topology changes occur only through authorized spawn operations recorded in the ledger; there is no mechanism for retroactive pointer manipulation.

\medskip
\noindent\textbf{Pre-commit hook.} The pre-commit hook is the enforcement mechanism that makes all other RSP constraints structural. At every spawn event, it verifies two conditions before the spawn is recorded: (a) $R(n_{i+1}) \subseteq R(n_i)$ and (b) $\mathrm{depth}(n_{i+1}) \leq D_{\mathrm{max}}$. Rejected spawns are logged; the Bounds Engine does not execute them; the protocol transitions to \texttt{ESCALATING} if the unmet condition persists.

Without the Fact Layer write protocol, immutability is a promise rather than a constraint. A mutable parent pointer enables retroactive chain manipulation. Without the pre-commit hook, the scope refinement constraint and depth bound are policy documentation — enforceable only until an actor with sufficient privilege decides they are inconvenient.

% -------------------------------------------------------
\subsection{Structural Minimality}
\label{sec:rs-minimality}

The three-layer mapping is minimal and sufficient. Removing any layer causes the corresponding guarantee to fail. Removing the Purpose Layer eliminates the exclusive human terminus and the exclusive source of spawn authorization policy. Removing the Bounds Engine eliminates the constrained execution evaluation layer. Removing the Fact Layer eliminates both the immutable parent pointer and the forensic ledger. No subset of layers achieves the same guarantees.

% -------------------------------------------------------
% SECTION 7
% -------------------------------------------------------
\section{Correctness Properties}
\label{sec:rs-correctness}

This section establishes three correctness properties for the Recursive Spawning Protocol: drift prevention, chain integrity, and termination. Each is stated formally, proved sketchedly, and connected to the BX3 layer that enforces it.

% -------------------------------------------------------
\subsection{Drift Prevention}
\label{sec:rs-drift-prevention}

\medskip
\noindent\textbf{Theorem (Drift Prevention).} Let $C = \langle n_0, n_1, \ldots, n_k \rangle$ be a Recursive Spawning chain rooted at $n_0$ with human anchor $h(n_0) = h(n_k)$. Then for all $i$ ($0 \leq i \leq k$), the operating scope $R(n_i)$ is a descendant refinement of the intent of $h(n_0)$.

\medskip
\noindent\textbf{Proof sketch.} The proof proceeds by structural induction on chain depth, where each inductive step is enforced at runtime by the Fact Layer pre-commit hook.

\medskip
\textit{Base case ($i=0$).} $R(n_0)$ is defined by $h(n_0)$ at provisioning. By the Purpose Layer's definition of authorized operating scope, $R(n_0)$ is a descendant refinement of $h(n_0)$'s intent by construction. No drift exists at the root.

\medskip
\textit{Inductive step.} Assume $R(n_i)$ is a descendant refinement of $h(n_0)$'s intent for some $i$ with $0 \leq i < k$. For $n_i$ to spawn $n_{i+1}$, the Fact Layer pre-commit hook verifies:

\begin{enumerate}
    \item[(a)] $R(n_{i+1}) \subseteq R(n_i)$ — the scope refinement constraint.
    \item[(b)] $\mathrm{depth}(n_{i+1}) \leq D_{\mathrm{max}}$ — the depth constraint.
\end{enumerate}

The Fact Layer rejects any spawn violating (a) or (b). Therefore $R(n_{i+1}) \subseteq R(n_i)$ holds for the child by structural enforcement. By the inductive hypothesis, $R(n_i)$ is a descendant refinement of $h(n_0)$'s intent. By transitivity of subset, $R(n_{i+1})$ is also a descendant refinement of $h(n_0)$'s intent. Hence $n_{i+1}$ does not drift. $\square$

\medskip
The drift prevention guarantee rests on the conjunction of three architectural facts:

\begin{enumerate}
    \item[(1)] \textbf{Immutable parent pointer.} $\alpha(n)$ is established once at provisioning and is thereafter immutable under the Fact Layer write protocol. Scope inheritance is traceable to a unique authorized lineage.
    \item[(2)] \textbf{Forensic ledger.} Every scope assignment is recorded in the append-only, hash-chained ledger. The inductive proof is empirically verifiable: each step of the refinement chain is a recorded, tamper-evident ledger entry.
    \item[(3)] \textbf{Chain terminates at Purpose Layer.} The exclusive human terminus closes the inductive argument at a finite bound. There is no machine-actor alternative terminus.
\end{enumerate}

If any of these three conditions were absent, drift prevention would not be structurally guaranteed. With all three, drift is architecturally impossible.

% -------------------------------------------------------
\subsection{Chain Integrity}
\label{sec:rs-chain-integrity}

\medskip
\noindent\textbf{Theorem (Chain Integrity).} For any Recursive Spawning chain $C = \langle n_0, n_1, \ldots, n_k \rangle$, the parent pointer sequence $\alpha(n_i) = n_{i-1}$ holds for all $1 \leq i \leq k$, and no node configuration in $C$ has been modified after provisioning.

\medskip
\noindent\textbf{Proof sketch.} The proof has two components.

\medskip
\textit{Topological integrity.} Each spawn event $n_{i-1} \xrightarrow{\mathrm{spawn}} n_i$ establishes $\alpha(n_i) = n_{i-1}$ as an immutable Fact Layer property at provisioning time. The Fact Layer write protocol rejects any operation attempting to modify $\alpha(n_i)$ after provisioning. The chain is constructed exclusively through authorized spawn events; no parent pointer is modified after establishment. This eliminates the re-parenting drift pathway.

\medskip
\textit{Node immutability.} Each node $n_i$ is provisioned with a fixed configuration $(R(n_i), P_{\mathrm{spawn}}(n_i), h(n_i))$ established by the Purpose Layer and recorded in the Fact Layer authorization ledger. The Fact Layer enforces immutability of these parameters after provisioning. The write protocol rejects any modification request from any actor. A \texttt{RESOLVE} directive is recorded as a new ledger entry (a state transition), not as an alteration of the provisioned configuration.

The forensic ledger's hash-chaining provides evidentially for chain integrity: any attempt to insert a falsified spawn event or reorder existing entries breaks hash continuity and is detectable. Chain integrity is not merely a property of the current chain state — it is a property of the entire historical record, enforceable at any future time. $\square$

% -------------------------------------------------------
\subsection{Termination}
\label{sec:rs-termination}

\medskip
\noindent\textbf{Theorem (Termination).} Every Recursive Spawning chain $C = \langle n_0, n_1, \ldots, n_k \rangle$ terminates — either at a resolution state or at the escalation threshold — within a finite number of spawn steps bounded by $\max(D_{\mathrm{max}}, D_{\mathrm{ESCALATE}})$.

\medskip
\noindent\textbf{Proof sketch.} The proof establishes two lemmas.

\medskip
\textit{Lemma 1 (Depth-bounded growth).} The chain grows only through spawn events. By the Bounds Engine depth enforcement and the Fact Layer pre-commit hook, any spawn event producing $|C| \geq D_{\mathrm{max}}$ is rejected. Therefore the chain cannot exceed depth $D_{\mathrm{max}} - 1$ through autonomous spawn. Rejected spawns are logged; the Bounds Engine transitions to \texttt{ESCALATING}.

\medskip
\textit{Lemma 2 (Escalation-triggered halt).} If the chain reaches $D_{\mathrm{ESCALATE}}$ before resolving the exception condition, the protocol enters \texttt{ESCALATING} state and suspends all autonomous spawn activity. The Bounds Engine awaits a directive from $h(n_k)$ — the human anchor. Since $h(n_k) = h(n_0)$ is a human principal, the directive arrives in finite time. \texttt{REJECT} terminates the chain definitively; \texttt{RESOLVE} redirects behavior and records a new ledger entry; \texttt{ACK} permits continued operation.

\medskip
\textit{Theorem conclusion.} Combining Lemma 1 and Lemma 2: the chain grows autonomously at most $D_{\mathrm{max}} - 1$ spawn steps. If it reaches $D_{\mathrm{ESCALATE}}$ before resolution, it halts and escalates. In all execution paths, the chain reaches a terminal state within $\max(D_{\mathrm{max}}, D_{\mathrm{ESCALATE}})$ spawn steps. Since both $D_{\mathrm{max}}$ and $D_{\mathrm{ESCALATE}}$ are finite integers defined at provisioning, the chain terminates in finite time. $\square$

The termination guarantee also precludes infinite spawn-escalate cycles. When $D_{\mathrm{max}}$ is reached, the Fact Layer blocks further spawn and triggers mandatory escalation. There is no mechanism by which the chain can circumvent this bound — the Fact Layer pre-commit hook is structural, not configurable at runtime by any machine actor.

% -------------------------------------------------------
\subsection{Collective Guarantee}
\label{sec:rs-collective-guarantee}

Drift prevention ensures scope refinement is the only authorized direction of chain growth. Chain integrity ensures the chain topology is stable and tamper-evident: the parent pointer sequence is immutable, node configurations are unmodifiable after provisioning, and the forensic ledger's hash-chaining makes any tampering detectable. Termination ensures the chain cannot grow indefinitely and cannot sustain an infinite spawn-escalate cycle: the depth bound is enforced structurally, and the escalation path terminates in human judgment, never in autonomous resolution.

Together, Theorems~\ref{sec:rs-drift-prevention}, \ref{sec:rs-chain-integrity}, and~\ref{sec:rs-termination} establish that a Recursive Spawning chain governed by the BX3 Framework is a provably accountable, integrity-preserving, and finite lineage of agents — suitable for deployment in high-stakes operational contexts where auditability and correctness are non-negotiable requirements.